\section{My team}

\textbf{2a. Structure, and whether we discussed it.} At our week-12 kickoff we drew up a division of labour on a whiteboard rather than in a document: Robin owned the mission state machine, Demetro owned aerodynamics and FDR presentation, Edward owned Cube wiring and bench debugging, our project manager owned budget and external comms, and I was nominally the CV lead. What we never explicitly discussed was the \emph{coupling} between those slots --- how my CV pipeline would hand data to Robin's state machine, how Demetro's airframe numbers would constrain my altitude assumptions, or how Edward's serial stack would reach my vision loop. I now recognise that my preference for a quick whiteboard allocation over a written interface contract was not efficiency; it was the same avoidance I show in personal planning --- a reluctance to spend an hour making the implicit explicit because it feels like bureaucracy. By week 17 the cost of that reluctance was visible: I was rewriting main.py modules that nobody else could safely touch, and Robin was writing a parallel flight script because integrating with mine was opaque. The uncomfortable realisation was that our ``structure'' had never actually been agreed --- it had been \emph{assumed}, and assumptions degrade silently.

\textbf{2b. What impacted our outputs most --- two conflicts, honestly told.} The first was the \textbf{11-state FSM disagreement} in week 17. I had grown main.py into an eleven-state machine (\texttt{INIT, CONNECTING, ARMING, TAKEOFF, SEARCH, CENTERING, DESCENDING, VERIFY, APPROACH, LANDING, DONE}) because each transition had caught a real bug in simulation. Robin read it and pushed back: he said it was unreadable and he would write his own shorter version. My first reaction was defensive --- I had the commit history and the pytest suite to prove every state earned its place. What I did instead of arguing was refactor: I split main.py from 1411 to 754 lines across six modules (\texttt{safety\_guard}, \texttt{navigation}, \texttt{state\_handlers}, \texttt{geofence}, \texttt{commands}, \texttt{telemetry}), wrote 146 pytest tests, and then built \texttt{passive\_watch\_2.py} with a \texttt{--robin} flag and an HTTP \texttt{cv\_mode} endpoint so Robin's script could poll vision without importing any of my code. The outcome was that Robin kept his short state machine and I kept mine, and the two ran in parallel on the Pi with no shared state --- documented in \texttt{docs/DESIGN\_DECISIONS.md:107-109}. I now see that what looked like a disagreement about code quality was actually a disagreement about \emph{ownership boundaries}, and the correct move was not to win the argument but to widen the interface so the argument became moot.

The second conflict was quieter but more personally difficult: the \textbf{workload-and-comms tension with our project manager} in late week 20, when three flight days had been cancelled and roughly eight weeks of planned hardware work had yielded zero seconds of real flight data. I felt the team was drifting into silent resignation and the PM was defaulting to ``wait for better weather.'' Instead of confronting him directly, which I have historically done badly, I drafted a complaint letter --- the whole team's grievance, in measured factual prose, with an appended chat history and a deliberate refusal to assign personal blame (\texttt{docs/COMPLAINT\_LETTER.md}). I iterated the tone six times (commits \texttt{26e3255, f3401e4, d946838, 5af87d0, 847518e, 4b0fd91}) before sharing it. The PM read it, disagreed with one paragraph, and we rewrote that paragraph together; the letter went out with all five names. The mediating action was \emph{externalising an internal conflict} --- converting a brewing interpersonal frustration into a shared document with a shared audience. I had not previously noticed in myself that I can only handle confrontation when it is mediated by writing. That is a pattern I now see across my personal life.

\textbf{2c. What I would change.} In a future project I would run a 90-minute \emph{interface workshop} in week 1, not week 12, and its output would be a written document naming every handoff (who calls what, with which data, at which rate) --- not a Gantt chart, not a Slack channel, a single shared file with signatures. I would also schedule a short \emph{reflection round} every two weeks --- not status, which we already had, but a round where each person names one thing that is degrading silently. The falsifiability test for whether this is working: I should be able to point, at the end of the project, to at least one conflict that was \emph{surfaced} in a reflection round before it became a rewrite. The deeper lesson --- the one I am still sitting with --- is that I came into this project with a Ukrainian engineering reflex of ``one person delivers when the system fails,'' and while that reflex produced the simulation framework and 820 commits, it also built a castle while the town starved. Next time I want to build fewer walls and more bridges, and the way I will know I have done it is if a teammate, not me, owns the integration point at the end.
